% ===================================================================
%  Physik: Schaltpläne mit circuitikz (deutsche Symbole: Option european)
%  Seite 1: Reihen- und Parallelschaltung
%  Seite 2: Gruppe A (Lampen) und Gruppe B (Widerstände) aus EINER Vorlage
%  Kompilieren: pdflatex vorlage_schaltplaene.tex
% ===================================================================
\documentclass[border=3mm]{standalone}
\usepackage[T1]{fontenc}
\usepackage{lmodern}
\usepackage[ngerman]{babel}
\usepackage{amsmath}
\usepackage[european]{circuitikz}

% Eine Schaltung mit Platzhalter-Bauteil: was "bauteil" ist, legt man davor fest
\newcommand\Schaltung[1]{%
  \begin{circuitikz}
    \draw (0,0) to[battery1, l=$U$] (0,3)
      to[bauteil, l=$X_1$] (3,3)
      to[bauteil, l=$X_2$] (3,0) -- (0,0);
    \node[font=\sffamily\bfseries] at (1.5,1.5) {#1};
  \end{circuitikz}}

\begin{document}
\begin{minipage}{14cm}
  \begin{circuitikz}   % Reihenschaltung mit Strommesser
    \draw (0,0) to[battery1, l=$U$] (0,3)
                to[lamp, l=$L_1$] (3,3)
                to[lamp, l=$L_2$] (3,0)
                to[rmeter, t=A] (0,0);
  \end{circuitikz}\hspace{2cm}
  \begin{circuitikz}   % Parallelschaltung mit Verzweigungspunkten
    \draw (0,0) to[battery1, l=$U$] (0,3)
                -- (2,3) to[lamp, l=$L_1$] (2,0)
                -- (0,0);
    \draw (2,3) -- (4,3) to[lamp, l=$L_2$] (4,0)
                -- (2,0);
    \draw (2,3) node[circ]{} (2,0) node[circ]{};
  \end{circuitikz}

  \bigskip
  \tikzset{bauteil/.style={lamp}}\Schaltung{A}\hspace{2cm}
  \tikzset{bauteil/.style={R}}\Schaltung{B}
\end{minipage}
\end{document}
